% ============================================================
% FEE3411 Control Systems A -- Tutorial sheet, Week 03
% Topic: Transfer functions, poles/zeros, and the s-plane
%
% ONE source, TWO outputs. The solutions toggle is driven by \jobname,
% so the two PDFs cannot drift apart:
%
%   latexmk -pdf -jobname=week-03-tutorial-questions week-03-tutorial.tex
%   latexmk -pdf -jobname=week-03-tutorial-solutions week-03-tutorial.tex
%
% Needs: ../tex/preamble.tex (this file lives in Tutorials/).
% ============================================================
\documentclass[11pt,a4paper]{article}
\usepackage[margin=25mm]{geometry}
% ============================================================
% FEE3411 Control Systems A -- shared preamble
% Used by both the lecture notes (article) and the slides (beamer).
% ============================================================
\usepackage[T1]{fontenc}
% lmodern gives cleaner PDF fonts but is absent from minimal TeX installs.
% Load it only if present, so this preamble builds anywhere.
\IfFileExists{lmodern.sty}{\usepackage{lmodern}}{}
\usepackage{amsmath,amssymb,amsthm}
\usepackage{graphicx}
\usepackage{booktabs}
\usepackage{siunitx}
\usepackage{tikz}
\usepackage{circuitikz}
\usepackage{pgfplots}
\pgfplotsset{compat=1.18}
\usetikzlibrary{arrows.meta,positioning,calc,decorations.markings,shapes.geometric}
\usepackage[most]{tcolorbox}
\usepackage{xcolor}

% ---- Course colours -------------------------------------------------
\definecolor{fnavy}{HTML}{1F3864}
\definecolor{faccent}{HTML}{2E5E4E}
\definecolor{fflag}{HTML}{9C4221}

% ---- Block-diagram style --------------------------------------------
\tikzset{
  block/.style   = {draw, thick, rectangle, minimum height=2.6em,
                    minimum width=3.2em, align=center},
  sumnode/.style = {draw, thick, circle, inner sep=0pt, minimum size=1.6em},
  pickoff/.style = {fill, circle, inner sep=0pt, minimum size=3pt},
  sig/.style     = {-{Latex[length=2mm]}, thick},
}

% ---- Summing-junction signs -----------------------------------------
% \sumsign{<node>}{<angle>}{<+ or ->} places a sign just outside the circle.
\newcommand{\sumsign}[3]{\node[font=\small] at ($(#1)+(#2:0.5)$) {$#3$};}

% ---- Boxed environments ---------------------------------------------
\newtcolorbox{keyidea}[1][]{colback=fnavy!5,colframe=fnavy,
  fonttitle=\bfseries,title=Key idea,#1}
\newtcolorbox{workedex}[1][]{colback=faccent!5,colframe=faccent,
  fonttitle=\bfseries,title=Worked example,#1}
\newtcolorbox{pitfall}[1][]{colback=fflag!5,colframe=fflag,
  fonttitle=\bfseries,title=Common pitfall,#1}

% ---- Shorthands ------------------------------------------------------
\newcommand{\Lap}{\mathcal{L}}
\newcommand{\jw}{\mathrm{j}\omega}
\newcommand{\wn}{\omega_{n}}
\newcommand{\zt}{\zeta}
\newcommand{\TF}[1]{#1(s)}

% ---- Week title machinery -------------------------------------------
\usepackage{fancyhdr}
\newcommand{\thecoursecode}{FEE3411}
\newcommand{\thecoursename}{Control Systems A}
\newcommand{\wknum}{}
\newcommand{\wktopic}{}
\newcommand{\weeknumber}[1]{\renewcommand{\wknum}{#1}}
\newcommand{\weektopic}[1]{\renewcommand{\wktopic}{#1}}
\newcommand{\weektitle}{%
  \thispagestyle{plain}%
  \begin{center}
    {\large\sffamily\bfseries\color{fnavy}\thecoursecode: \thecoursename}\\[2pt]
    {\Huge\sffamily\bfseries Week \wknum}\\[4pt]
    {\LARGE\sffamily\color{fnavy}\wktopic}\\[6pt]
    \rule{\textwidth}{1.2pt}
  \end{center}\vspace{4pt}%
  \pagestyle{fancy}\fancyhf{}%
  \fancyhead[L]{\footnotesize\color{gray}\thecoursecode\ Week \wknum\ --- \wktopic}%
  \fancyhead[R]{\footnotesize\color{gray}\thepage}%
  \renewcommand{\headrulewidth}{0.4pt}%
}

% ---- Outcomes and reading boxes -------------------------------------
\newenvironment{outcomes}
  {\begin{tcolorbox}[colback=white,colframe=fnavy!60,
      fonttitle=\bfseries,title=By the end of this week you should be able to]
   \begin{itemize}[leftmargin=*]}
  {\end{itemize}\end{tcolorbox}}

\newtcolorbox{readingbox}{colback=gray!5,colframe=gray!50,
  fonttitle=\bfseries,title=Reading}

% enumitem must NOT be loaded under beamer: it overwrites beamer's itemize
% templates and the bullet markers silently disappear. The notes (article) get
% it; the slides (beamer) use beamer's own list machinery instead.
\makeatletter
\@ifclassloaded{beamer}{}{\usepackage{enumitem}}
\makeatother

\usepackage{environ}

% ---- solutions toggle, decided by the job name ----------------------
% NB: xstring's \IfSubStr does NOT work here -- \jobname yields catcode-12
% characters and the comparison silently fails, giving the questions version
% both times. The expl3 string test below is catcode-blind and does work.
\newif\ifsolutions
\ExplSyntaxOn
\tl_set:Nx \l_tmpa_tl { \jobname }
\str_if_in:NnTF \l_tmpa_tl { solutions } { \solutionstrue } { \solutionsfalse }
\ExplSyntaxOff

\newtcolorbox{solbox}[1][]{colback=faccent!4,colframe=faccent!70,
  fonttitle=\bfseries,title=Solution,breakable,#1}
\NewEnviron{solution}{\ifsolutions\begin{solbox}\BODY\end{solbox}\fi}

% ---- question tagging ------------------------------------------------
\newcommand{\inhour}{\hfill{\small\color{faccent}\itshape[in the hour]}}
\newcommand{\athome}{\hfill{\small\color{fflag}\itshape[homework]}}

\weeknumber{3}
\weektopic{Tutorial 3 --- Transfer Functions and the $s$-Plane}

\begin{document}
\weektitle

\begin{center}
\ifsolutions
  {\large\sffamily\bfseries\color{fflag}QUESTIONS AND SOLUTIONS}\\[2pt]
  {\small\itshape Instructor copy --- do not circulate before the tutorial.}
\else
  {\large\sffamily\bfseries\color{fnavy}QUESTION SHEET}
\fi
\end{center}

\bigskip
\noindent
This hour delivers on the two promises made at the end of the Week~3 lecture:
full transfer-function derivations, from physical law to $G(s)$, for an
electrical system (Question~1) and a mechanical system (Question~2); and more
practice with inverse Laplace transforms by partial fractions, extending the
notes to the two cases they describe but do not fully work through with
numbers --- a repeated pole (Question~3) and a right-half-plane zero
(Question~4) --- plus a second, distinct complex-pole example
(Question~5).

\smallskip
\noindent
Questions marked \textbf{\color{faccent}[in the hour]} are the ones we work
through together. Those marked \textbf{\color{fflag}[homework]} you do in your
own time --- they are not collected, but the technique is assumed from here
on, including in Assignment~2 (issued Week~4).

\smallskip
\noindent
\emph{Preparation:} Week~3 lecture notes, \S1--\S4.

% --------------------------------------------------------------------
\begin{readingbox}
\textbf{Quick reference --- carry these into every question below.}

\smallskip
\begin{minipage}[t]{0.48\textwidth}
\textbf{Transfer function, zero initial conditions}\\[2pt]
\footnotesize
\[
  G(s) = \frac{Y(s)}{U(s)}\bigg|_{\text{zero i.c.}}
\]
\[
  \dot y \leftrightarrow sY(s),\qquad
  \ddot y \leftrightarrow s^2Y(s)
\]
\textbf{Impedance shortcut (zero i.c.\ only)}
\[
  Z_R=R,\qquad Z_L=Ls,\qquad Z_C=\frac{1}{Cs}
\]
\textbf{d.c.\ gain (final value, unit step input)}
\[
  y(\infty)=\lim_{s\to0}sY(s)=G(0)
\]
(poles of $sY(s)$ all in the open left half plane)
\end{minipage}
\hfill
\begin{minipage}[t]{0.48\textwidth}
\textbf{Partial-fraction templates}\\[2pt]
\footnotesize
\emph{Distinct real pole:}
\[
  \frac{A}{s+a} \leftrightarrow Ae^{-at}
\]
\emph{Repeated pole, order 2:}
\[
  \frac{B}{s+a}+\frac{C}{(s+a)^2}
  \leftrightarrow
  Be^{-at}+Cte^{-at}
\]
\emph{Complex-conjugate pair, $(s+a)^2+\omega^2$:}
\[
  \frac{s+a}{(s+a)^2+\omega^2}\leftrightarrow e^{-at}\cos\omega t
\]
\[
  \frac{\omega}{(s+a)^2+\omega^2}\leftrightarrow e^{-at}\sin\omega t
\]
\end{minipage}
\end{readingbox}

% ====================================================================
\section*{Question 1 --- Series $RLC$ circuit: full derivation \inhour}

A series $RLC$ circuit is driven by a voltage source $v(t)$, and the output is
taken as the voltage across the capacitor, $v_C(t)$. The component values are
$R=\SI{7}{\ohm}$, $L=\SI{1}{\henry}$, $C=\SI{0.1}{\farad}$.

\begin{enumerate}[label=(\alph*)]
\item Write Kirchhoff's voltage law around the loop in terms of the loop
      current $i(t)$, and the defining relation $i(t)=C\,\dot v_C(t)$. Combine
      the two into a single second-order differential equation in $v_C(t)$
      and $v(t)$.
\item Take the Laplace transform of your equation with zero initial
      conditions and hence write down $G(s)=V_C(s)/V(s)$.
\item Factor the denominator to find the two poles. Which is dominant, and
      what is its time constant?
\item Find the d.c.\ gain $G(0)$ two ways: (i) directly from your transfer
      function, and (ii) by reasoning physically about the circuit once all
      transients have died away. Confirm the two agree.
\end{enumerate}

\begin{solution}
\textbf{(a)} KVL around the loop: $v = Ri + L\dot i + v_C$, with
$i=C\dot v_C$. Substituting,
\[
  v = RC\dot v_C + LC\ddot v_C + v_C
  \quad\Longleftrightarrow\quad
  LC\,\ddot v_C + RC\,\dot v_C + v_C = v .
\]

\medskip
\textbf{(b)} With zero initial conditions, $\ddot v_C\to s^2V_C(s)$,
$\dot v_C\to sV_C(s)$:
\[
  (LCs^2+RCs+1)V_C(s) = V(s)
  \quad\Longrightarrow\quad
  G(s)=\frac{V_C(s)}{V(s)}=\frac{1}{LCs^2+RCs+1}.
\]
Substituting $L=1$, $C=0.1$, $R=7$: $LC=0.1$, $RC=0.7$, so
\[
  G(s)=\frac{1}{0.1s^2+0.7s+1}=\frac{10}{s^2+7s+10}.
\]

\medskip
\textbf{(c)} Factoring, $s^2+7s+10=(s+2)(s+5)$, so the poles are
$s=-2$ and $s=-5$. The pole at $s=-2$ is dominant (closest to the origin,
decays slowest), with time constant $\tau=1/2=\SI{0.5}{\second}$; the pole
at $s=-5$ decays two and a half times faster ($\tau=0.2\,\si{\second}$) and
its transient is essentially gone before the dominant one is.

\medskip
\textbf{(d)} (i) $G(0)=10/10=1$. (ii) Once transients die away with a
constant (d.c.) source, no current flows into the capacitor
($i=C\dot v_C=0$ at steady state), so there is no drop across $R$ or $L$ and
$v_C=v$: the d.c.\ gain is exactly $1$, confirming (i).
\end{solution}

% ====================================================================
\section*{Question 2 --- Mass--spring--damper: full derivation \inhour}

A mass $M=\SI{1}{\kilogram}$ slides on a damped surface and is connected to a
wall by a spring, with damping coefficient $B=\SI{7}{\newton\second\per\metre}$
and spring constant $K=\SI{12}{\newton\per\metre}$. A force $f(t)$ is applied
to the mass, and the output is its displacement $x(t)$ from equilibrium.

\begin{enumerate}[label=(\alph*)]
\item Apply Newton's second law to write the equation of motion in $x(t)$
      and $f(t)$.
\item Take the Laplace transform with zero initial conditions and find
      $G(s)=X(s)/F(s)$.
\item Factor the denominator to find the two poles, and state which is
      dominant.
\item Recall the force--voltage analogy from Week~2 ($M\!\leftrightarrow\! L$,
      $B\!\leftrightarrow\! R$, $K\!\leftrightarrow\! 1/C$). What values of
      $R$, $L$, $C$ would give an electrical circuit with exactly this
      transfer function shape? Compare the resulting $C$ with Question~1's
      --- are the two systems' pole locations the same? Should they be?
\end{enumerate}

\begin{solution}
\textbf{(a)} Newton's second law, with spring and damper forces opposing the
motion: $M\ddot x = f - B\dot x - Kx$, i.e.
\[
  M\ddot x + B\dot x + Kx = f .
\]

\medskip
\textbf{(b)} With zero initial conditions,
\[
  (Ms^2+Bs+K)X(s)=F(s)
  \quad\Longrightarrow\quad
  G(s)=\frac{X(s)}{F(s)}=\frac{1}{Ms^2+Bs+K}
  =\frac{1}{s^2+7s+12}.
\]

\medskip
\textbf{(c)} Factoring, $s^2+7s+12=(s+3)(s+4)$: poles at $s=-3$ and $s=-4$.
The pole at $s=-3$ is dominant, time constant $\tau=1/3=\SI{0.333}{\second}$.

\medskip
\textbf{(d)} Matching $M\leftrightarrow L=\SI{1}{\henry}$,
$B\leftrightarrow R=\SI{7}{\ohm}$, $K\leftrightarrow 1/C\Rightarrow
C=1/12=\SI{0.0833}{\farad}$. This is close to, but not the same as,
Question~1's $C=\SI{0.1}{\farad}$ --- and correspondingly the pole locations
differ ($-3,-4$ here against $-2,-5$ there), even though both transfer
functions have the identical form $1/(s^2+7s+K')$ with the same middle
coefficient. The two systems are analogous in structure, not identical in
behaviour: the analogy carries over the \emph{equation}, not the specific
numbers, and a small change in the third coefficient is enough to move both
poles.
\end{solution}

\begin{pitfall}
Questions~1 and~2 both reduce to $G(s)=K'/(s^2+7s+K')$ with the \emph{same}
$R/L$ (or $B/M$) ratio of $7$, purely because these examples were chosen that
way. Do not read anything physical into that coincidence beyond ``both
systems happen to have the same relative damping term.'' The poles, the
d.c.\ gain, and the time constants are all set by the actual numbers in
\emph{each} system, and Question~1's poles ($-2,-5$) are not interchangeable
with Question~2's ($-3,-4$) just because the equations look alike.
\end{pitfall}

% ====================================================================
\section*{Question 3 --- Partial fractions with a repeated pole \inhour}

A system has transfer function
\[
  G(s)=\frac{8}{(s+2)^2}.
\]

\begin{enumerate}[label=(\alph*)]
\item Find the d.c.\ gain $G(0)$, and hence predict the final value of the
      unit-step response \emph{before} doing any partial-fraction work.
\item Form $Y(s)=G(s)/s$ and expand it in partial fractions of the form
      $A/s+B/(s+2)+C/(s+2)^2$. Find $A$, $B$, $C$.
\item Use the repeated-pole template from the quick-reference box to invert
      $Y(s)$ to $y(t)$. Which term is new compared with the distinct-pole
      case, and where does it come from mathematically (what did the second
      power in $(s+2)^2$ demand)?
\item Confirm $\lim_{t\to\infty}y(t)$ matches your prediction from (a).
\end{enumerate}

\begin{solution}
\textbf{(a)} $G(0)=8/4=2$, so we expect $y(\infty)=2$ (the poles of $sY(s)$,
namely $s=0$ and the double pole $s=-2$, are all in the closed left half
plane with only a simple pole at the origin, so the final value theorem
applies).

\medskip
\textbf{(b)}
\[
  Y(s)=\frac{8}{s(s+2)^2}=\frac{A}{s}+\frac{B}{s+2}+\frac{C}{(s+2)^2}.
\]
Cover-up for $A$ and $C$: $A=\left.\dfrac{8}{(s+2)^2}\right|_{s=0}=\dfrac{8}{4}=2$;
$C=\left.\dfrac{8}{s}\right|_{s=-2}=\dfrac{8}{-2}=-4$. For $B$, multiply
through by $s(s+2)^2$ and match the $s^2$ coefficient: $0=A+B\Rightarrow
B=-A=-2$. So
\[
  Y(s)=\frac{2}{s}-\frac{2}{s+2}-\frac{4}{(s+2)^2}.
\]

\medskip
\textbf{(c)} Using $B/(s+a)\leftrightarrow Be^{-at}$ and
$C/(s+a)^2\leftrightarrow Cte^{-at}$ with $a=2$:
\[
  y(t)=2-2e^{-2t}-4te^{-2t},\qquad t\ge0.
\]
The new term is $-4te^{-2t}$: it did not appear in Question~1 or~2, where
every pole was simple. A \emph{repeated} pole at $s=-2$ needs a second,
independent basis function alongside $e^{-2t}$ to match the extra degree of
freedom in the partial fraction, and $te^{-2t}$ is that function --- it is
still decaying at the same rate $e^{-2t}$, but it starts at zero and rises
before falling, so it briefly reinforces the decay rather than opposing it.

\medskip
\textbf{(d)} As $t\to\infty$, both $e^{-2t}\to0$ and $te^{-2t}\to0$ (the
exponential always beats the polynomial), leaving $y(\infty)=2$, exactly
matching (a).
\end{solution}

\begin{keyidea}
A repeated real pole at $s=-a$ is the borderline case between two distinct
real poles and a complex-conjugate pair: it is what you get from
$s^2+2as+a^2=(s+a)^2$, the exact boundary of $\zeta=1$ (critical damping) in
the standard second-order form $s^2+2\zeta\wn s+\wn^2$. That is why it needs
its own template ($te^{-at}$, not a second, distinct exponential): the two
roots have merged, and the system has ``used up'' its independence between
modes. You will meet this same boundary again when second-order response
specifications (overshoot, settling time) are built from $\zeta$ and $\wn$.
\end{keyidea}

% ====================================================================
\section*{Question 4 --- A right-half-plane zero \athome}

Two systems share the same two poles, $s=-2$ and $s=-6$, and the same
d.c.\ gain of $1$, but differ in their zero:
\[
  G_A(s)=\frac{12}{(s+2)(s+6)}
  \qquad\text{(no zero)},
\]
\[
  G_B(s)=\frac{K(s-4)}{(s+2)(s+6)}
  \qquad\text{(zero at }s=+4\text{, right half plane)}.
\]

\begin{enumerate}[label=(\alph*)]
\item Find the value of $K$ that gives $G_B(0)=1$, matching $G_A$'s d.c.\
      gain. (Be careful with the sign.)
\item Expand $Y_A(s)=G_A(s)/s$ in partial fractions and invert it to find
      $y_A(t)$.
\item Do the same for $Y_B(s)=G_B(s)/s$ to find $y_B(t)$.
\item Compute $\dot y_A(0)$ and $\dot y_B(0)$ directly from your time-domain
      expressions. What is qualitatively different about the two responses in
      the first instant, and which feature of $G_B$ is responsible?
\item Using your expression for $y_B(t)$ (or a calculator/plot), estimate how
      far below zero $y_B(t)$ dips, and roughly when. Does $y_B(t)$ ever
      become unstable?
\end{enumerate}

\begin{solution}
\textbf{(a)} $G_B(0)=K(-4)/[(2)(6)]=-K/3$. Setting this to $1$ gives
$K=-3$, so
\[
  G_B(s)=\frac{-3(s-4)}{(s+2)(s+6)}=\frac{3(4-s)}{(s+2)(s+6)}.
\]
The sign matters: a right-half-plane zero at $+4$ needs a \emph{negative}
leading coefficient to keep the d.c.\ gain positive --- it is easy to guess
$K=3$ by pattern-matching the numerator to $(4-s)$ without checking, and get
the wrong answer.

\medskip
\textbf{(b)} $Y_A(s)=\dfrac{12}{s(s+2)(s+6)}=\dfrac{A}{s}+\dfrac{B}{s+2}+\dfrac{C}{s+6}$.
Cover-up: $A=12/12=1$; $B=12/[(-2)(4)]=-1.5$; $C=12/[(-6)(-4)]=0.5$. So
\[
  y_A(t)=1-1.5e^{-2t}+0.5e^{-6t}.
\]

\medskip
\textbf{(c)} $Y_B(s)=\dfrac{3(4-s)}{s(s+2)(s+6)}=\dfrac{A}{s}+\dfrac{B}{s+2}+\dfrac{C}{s+6}$.
$A=G_B(0)=1$; $B=\left.\dfrac{3(4-s)}{s(s+6)}\right|_{s=-2}=\dfrac{3(6)}{(-2)(4)}=-2.25$;
$C=\left.\dfrac{3(4-s)}{s(s+2)}\right|_{s=-6}=\dfrac{3(10)}{(-6)(-4)}=1.25$. So
\[
  y_B(t)=1-2.25e^{-2t}+1.25e^{-6t}.
\]

\medskip
\textbf{(d)} $\dot y_A(t)=3e^{-2t}-3e^{-6t}\Rightarrow\dot y_A(0)=0$.
$\dot y_B(t)=4.5e^{-2t}-7.5e^{-6t}\Rightarrow\dot y_B(0)=4.5-7.5=-3$.
System~$A$ leaves the origin with zero slope (it has relative degree~2, like
Question~3); system~$B$ leaves with a \emph{negative} slope even though it
is heading towards $+1$ --- it initially moves the wrong way. This is exactly
the right-half-plane zero at work: it is the only structural difference
between $G_A$ and $G_B$, and it is what forces the initial dip.

\medskip
\textbf{(e)} Evaluating $y_B(t)$ numerically, the minimum is
$y_B\approx-0.162$ at $t\approx0.13\,\si{\second}$ --- an undershoot of about
$16\%$ of the final value, arriving quickly and then recovering. $y_B(t)$
is not unstable: both poles ($-2,-6$) are in the left half plane, so
$y_B(t)\to1$ as $t\to\infty$ exactly as $y_A(t)$ does. The zero changes the
\emph{shape} of the transient, never whether the system settles.
\end{solution}

% ====================================================================
\section*{Question 5 --- Partial fractions with a complex-conjugate pair \athome}

A system has transfer function
\[
  G(s)=\frac{20}{s^2+4s+20}.
\]

\begin{enumerate}[label=(\alph*)]
\item Complete the square in the denominator to write it as
      $(s+a)^2+\omega^2$, and hence identify the poles, $\wn$, and $\zt$.
\item Find the d.c.\ gain $G(0)$.
\item Form $Y(s)=G(s)/s$ and expand it as $A/s+(Bs+C)/(s^2+4s+20)$. Find
      $A$, $B$, $C$.
\item Rewrite the second term using the same completed-square denominator
      from (a), split it into a $\cos$-generating part and a
      $\sin$-generating part (matching the templates in the quick-reference
      box), and hence write $y(t)$ in closed form.
\item Check: does your $y(t)$ give $y(0)=0$ and $y(\infty)=G(0)$?
\end{enumerate}

\begin{solution}
\textbf{(a)} $s^2+4s+20=(s+2)^2+16$, so the poles are $s=-2\pm\mathrm{j}4$.
Matching $(s+2)^2+16$ to the standard form $(s+\zt\wn)^2+\wn^2(1-\zt^2)$
gives $\wn^2=20\Rightarrow\wn=\sqrt{20}=2\sqrt5\approx\SI{4.47}{\radian\per\second}$,
and $\zt\wn=2\Rightarrow\zt=2/4.47\approx0.447$.

\medskip
\textbf{(b)} $G(0)=20/20=1$.

\medskip
\textbf{(c)}
\[
  Y(s)=\frac{20}{s(s^2+4s+20)}=\frac{A}{s}+\frac{Bs+C}{s^2+4s+20}.
\]
At $s=0$: $20=20A\Rightarrow A=1$. Matching the $s^2$ coefficient:
$0=A+B\Rightarrow B=-1$. Matching the $s^1$ coefficient:
$0=4A+C\Rightarrow C=-4$. So
\[
  Y(s)=\frac{1}{s}-\frac{s+4}{s^2+4s+20}=\frac{1}{s}-\frac{s+4}{(s+2)^2+16}.
\]

\medskip
\textbf{(d)} Split $s+4=(s+2)+2$:
\[
  \frac{s+4}{(s+2)^2+16}=\underbrace{\frac{s+2}{(s+2)^2+16}}_{\to\,e^{-2t}\cos4t}
  +\underbrace{\frac{2}{(s+2)^2+16}}_{=\frac12\cdot\frac{4}{(s+2)^2+16}\,\to\,\frac12e^{-2t}\sin4t}.
\]
Hence
\[
  \boxed{\;y(t)=1-e^{-2t}\left(\cos4t+\tfrac12\sin4t\right)\;},\qquad t\ge0.
\]

\medskip
\textbf{(e)} At $t=0$: $y(0)=1-(1)(1+0)=0$, confirming the step response
starts from rest. As $t\to\infty$, $e^{-2t}\to0$ (the oscillation is damped
by the same real part $-2$ that appears in both poles), leaving $y(\infty)=1$,
matching $G(0)$ from (b).
\end{solution}

% ====================================================================
\vfill
\begin{center}
\rule{0.5\textwidth}{0.4pt}\\[4pt]
{\small\itshape
\ifsolutions
  End of solutions. Week~4: block diagrams and signal-flow graphs.
\else
  Solutions to the homework questions are released with the Week~4 sheet.\\
  Next week: block diagrams and signal-flow graphs.
\fi}
\end{center}

\end{document}
